% ================================================================================
% PREAMBULO COMPARTIDO — Manuales de GAMMA
%
% Lo usan los tres manuales (usuario, tecnico e instalacion) para que compartan
% tipografia, paleta y recuadros. Se incluye con % ================================================================================
% PREAMBULO COMPARTIDO — Manuales de GAMMA
%
% Lo usan los tres manuales (usuario, tecnico e instalacion) para que compartan
% tipografia, paleta y recuadros. Se incluye con % ================================================================================
% PREAMBULO COMPARTIDO — Manuales de GAMMA
%
% Lo usan los tres manuales (usuario, tecnico e instalacion) para que compartan
% tipografia, paleta y recuadros. Se incluye con % ================================================================================
% PREAMBULO COMPARTIDO — Manuales de GAMMA
%
% Lo usan los tres manuales (usuario, tecnico e instalacion) para que compartan
% tipografia, paleta y recuadros. Se incluye con \input{../preambulo} desde la
% carpeta de cada manual.
%
% La paleta sale de las variables de frontend/src/style.css: la interfaz es
% enteramente acromatica, y los manuales tampoco introducen color.
% ================================================================================
\documentclass[11pt, letterpaper]{article}

\usepackage[spanish,mexico]{babel}
\usepackage[utf8]{inputenc}
\usepackage[T1]{fontenc}

% Tipografia: Inter, la mas cercana en TeX Live a la Geist que usa la
% aplicacion. Sin serifas en todo el documento, igual que la interfaz.
\usepackage[sfdefault]{inter}
\usepackage[scaled=0.86]{beramono}

\usepackage[top=2.4cm, bottom=2.4cm, left=2.6cm, right=2.6cm]{geometry}
\usepackage{graphicx}
\usepackage{xcolor}
\usepackage{enumitem}
\usepackage{fancyhdr}
\usepackage{titlesec}
\usepackage{tcolorbox}
% `enhanced` y `borderline west` viven en skins; `breakable` en su propia
% libreria. Sin declararlas, tcolorbox ignora esas claves en silencio y los
% recuadros pierden la barra lateral que los distingue.
\tcbuselibrary{skins,breakable}
\usepackage{booktabs}
\usepackage[hidelinks]{hyperref}

% ── Paleta ──────────────────────────────────────────────────────────────────
% Tomada de las variables de src/style.css. La interfaz es enteramente
% acromatica —todos sus colores son oklch(L 0 0)— asi que el manual tampoco
% introduce color: solo grises, con el texto como unico elemento de contraste.
\definecolor{fg}{HTML}{0A0A0A}        % --foreground
\definecolor{fgsoft}{HTML}{171717}    % --primary
\definecolor{muted}{HTML}{737373}     % --muted-foreground
\definecolor{surface}{HTML}{F5F5F5}   % --muted / --secondary
\definecolor{surfacealt}{HTML}{FAFAFA}% --sidebar
\definecolor{line}{HTML}{E5E5E5}      % --border

\titleformat{\section}{\LARGE\bfseries\color{fg}}{\thesection}{0.7em}{}
\titlespacing*{\section}{0pt}{1.6em}{0.7em}
\titleformat{\subsection}{\large\bfseries\color{fgsoft}}{\thesubsection}{0.6em}{}
\titlespacing*{\subsection}{0pt}{1.3em}{0.4em}
\titleformat{\subsubsection}[runin]{\bfseries\color{fgsoft}}{}{0pt}{}[.]

\pagestyle{fancy}
\fancyhf{}
\lhead{\footnotesize\color{muted}Manual de usuario}
\rhead{\footnotesize\color{muted}\thepage}
\renewcommand{\headrulewidth}{0.4pt}
\renewcommand{\headrule}{\color{line}\hrule height 0.4pt}

\setlength{\parindent}{0pt}
\setlength{\parskip}{0.62em}

% Enlaces y referencias en el mismo gris del cuerpo: sin subrayados ni azules.
\hypersetup{colorlinks=false}

% ── Recuadros ───────────────────────────────────────────────────────────────
% Se distinguen por una barra lateral y el peso del titulo, no por color.
\newtcolorbox{nota}[1][Nota]{
  enhanced, breakable, colback=surfacealt, colframe=surfacealt,
  borderline west={2pt}{0pt}{line},
  coltitle=fg, fonttitle=\bfseries\footnotesize\sffamily,
  title=#1, boxrule=0pt, arc=0pt, outer arc=0pt,
  left=10pt, right=10pt, top=7pt, bottom=7pt, toptitle=1pt, bottomtitle=3pt}

\newtcolorbox{aviso}[1][Importante]{
  enhanced, breakable, colback=surface, colframe=surface,
  borderline west={2pt}{0pt}{fg},
  coltitle=fg, fonttitle=\bfseries\footnotesize\sffamily,
  title=#1, boxrule=0pt, arc=0pt, outer arc=0pt,
  left=10pt, right=10pt, top=7pt, bottom=7pt, toptitle=1pt, bottomtitle=3pt}

% ── Captura de pantalla con marcador de posicion ────────────────────────────
% Uso: \captura{nombre-archivo}{Pie de la figura}
% Si manual/capturas/<nombre>.png existe, la inserta. Si no, dibuja el hueco.
\newcommand{\captura}[2]{%
  \begin{figure}[htbp]
  \centering
  \IfFileExists{capturas/#1.png}{%
    % Ancho completo, pero acotado en alto: una captura de ventana completa
    % puede ser muy alta y desbordaria la pagina. keepaspectratio evita que se
    % deforme cuando manda la restriccion de altura.
    {\color{line}\fboxrule=0.6pt\fboxsep=0pt\fbox{%
      \includegraphics[width=0.94\textwidth, height=0.62\textheight, keepaspectratio]{capturas/#1.png}}}%
  }{%
    \fcolorbox{line}{surfacealt}{%
      \begin{minipage}[c][3.1cm][c]{0.93\textwidth}
        \centering
        {\footnotesize\bfseries\color{muted}CAPTURA PENDIENTE}\\[5pt]
        {\ttfamily\footnotesize\color{fgsoft} capturas/#1.png}\\[4pt]
        {\footnotesize\color{muted} #2}
      \end{minipage}}%
  }
  \caption{#2}
  \label{fig:#1}
  \end{figure}%
}

% Pies de figura discretos
\usepackage{caption}
\captionsetup{font={small,color=muted}, labelfont={bf,color=muted}, skip=6pt}

 desde la
% carpeta de cada manual.
%
% La paleta sale de las variables de frontend/src/style.css: la interfaz es
% enteramente acromatica, y los manuales tampoco introducen color.
% ================================================================================
\documentclass[11pt, letterpaper]{article}

\usepackage[spanish,mexico]{babel}
\usepackage[utf8]{inputenc}
\usepackage[T1]{fontenc}

% Tipografia: Inter, la mas cercana en TeX Live a la Geist que usa la
% aplicacion. Sin serifas en todo el documento, igual que la interfaz.
\usepackage[sfdefault]{inter}
\usepackage[scaled=0.86]{beramono}

\usepackage[top=2.4cm, bottom=2.4cm, left=2.6cm, right=2.6cm]{geometry}
\usepackage{graphicx}
\usepackage{xcolor}
\usepackage{enumitem}
\usepackage{fancyhdr}
\usepackage{titlesec}
\usepackage{tcolorbox}
% `enhanced` y `borderline west` viven en skins; `breakable` en su propia
% libreria. Sin declararlas, tcolorbox ignora esas claves en silencio y los
% recuadros pierden la barra lateral que los distingue.
\tcbuselibrary{skins,breakable}
\usepackage{booktabs}
\usepackage[hidelinks]{hyperref}

% ── Paleta ──────────────────────────────────────────────────────────────────
% Tomada de las variables de src/style.css. La interfaz es enteramente
% acromatica —todos sus colores son oklch(L 0 0)— asi que el manual tampoco
% introduce color: solo grises, con el texto como unico elemento de contraste.
\definecolor{fg}{HTML}{0A0A0A}        % --foreground
\definecolor{fgsoft}{HTML}{171717}    % --primary
\definecolor{muted}{HTML}{737373}     % --muted-foreground
\definecolor{surface}{HTML}{F5F5F5}   % --muted / --secondary
\definecolor{surfacealt}{HTML}{FAFAFA}% --sidebar
\definecolor{line}{HTML}{E5E5E5}      % --border

\titleformat{\section}{\LARGE\bfseries\color{fg}}{\thesection}{0.7em}{}
\titlespacing*{\section}{0pt}{1.6em}{0.7em}
\titleformat{\subsection}{\large\bfseries\color{fgsoft}}{\thesubsection}{0.6em}{}
\titlespacing*{\subsection}{0pt}{1.3em}{0.4em}
\titleformat{\subsubsection}[runin]{\bfseries\color{fgsoft}}{}{0pt}{}[.]

\pagestyle{fancy}
\fancyhf{}
\lhead{\footnotesize\color{muted}Manual de usuario}
\rhead{\footnotesize\color{muted}\thepage}
\renewcommand{\headrulewidth}{0.4pt}
\renewcommand{\headrule}{\color{line}\hrule height 0.4pt}

\setlength{\parindent}{0pt}
\setlength{\parskip}{0.62em}

% Enlaces y referencias en el mismo gris del cuerpo: sin subrayados ni azules.
\hypersetup{colorlinks=false}

% ── Recuadros ───────────────────────────────────────────────────────────────
% Se distinguen por una barra lateral y el peso del titulo, no por color.
\newtcolorbox{nota}[1][Nota]{
  enhanced, breakable, colback=surfacealt, colframe=surfacealt,
  borderline west={2pt}{0pt}{line},
  coltitle=fg, fonttitle=\bfseries\footnotesize\sffamily,
  title=#1, boxrule=0pt, arc=0pt, outer arc=0pt,
  left=10pt, right=10pt, top=7pt, bottom=7pt, toptitle=1pt, bottomtitle=3pt}

\newtcolorbox{aviso}[1][Importante]{
  enhanced, breakable, colback=surface, colframe=surface,
  borderline west={2pt}{0pt}{fg},
  coltitle=fg, fonttitle=\bfseries\footnotesize\sffamily,
  title=#1, boxrule=0pt, arc=0pt, outer arc=0pt,
  left=10pt, right=10pt, top=7pt, bottom=7pt, toptitle=1pt, bottomtitle=3pt}

% ── Captura de pantalla con marcador de posicion ────────────────────────────
% Uso: \captura{nombre-archivo}{Pie de la figura}
% Si manual/capturas/<nombre>.png existe, la inserta. Si no, dibuja el hueco.
\newcommand{\captura}[2]{%
  \begin{figure}[htbp]
  \centering
  \IfFileExists{capturas/#1.png}{%
    % Ancho completo, pero acotado en alto: una captura de ventana completa
    % puede ser muy alta y desbordaria la pagina. keepaspectratio evita que se
    % deforme cuando manda la restriccion de altura.
    {\color{line}\fboxrule=0.6pt\fboxsep=0pt\fbox{%
      \includegraphics[width=0.94\textwidth, height=0.62\textheight, keepaspectratio]{capturas/#1.png}}}%
  }{%
    \fcolorbox{line}{surfacealt}{%
      \begin{minipage}[c][3.1cm][c]{0.93\textwidth}
        \centering
        {\footnotesize\bfseries\color{muted}CAPTURA PENDIENTE}\\[5pt]
        {\ttfamily\footnotesize\color{fgsoft} capturas/#1.png}\\[4pt]
        {\footnotesize\color{muted} #2}
      \end{minipage}}%
  }
  \caption{#2}
  \label{fig:#1}
  \end{figure}%
}

% Pies de figura discretos
\usepackage{caption}
\captionsetup{font={small,color=muted}, labelfont={bf,color=muted}, skip=6pt}

 desde la
% carpeta de cada manual.
%
% La paleta sale de las variables de frontend/src/style.css: la interfaz es
% enteramente acromatica, y los manuales tampoco introducen color.
% ================================================================================
\documentclass[11pt, letterpaper]{article}

\usepackage[spanish,mexico]{babel}
\usepackage[utf8]{inputenc}
\usepackage[T1]{fontenc}

% Tipografia: Inter, la mas cercana en TeX Live a la Geist que usa la
% aplicacion. Sin serifas en todo el documento, igual que la interfaz.
\usepackage[sfdefault]{inter}
\usepackage[scaled=0.86]{beramono}

\usepackage[top=2.4cm, bottom=2.4cm, left=2.6cm, right=2.6cm]{geometry}
\usepackage{graphicx}
\usepackage{xcolor}
\usepackage{enumitem}
\usepackage{fancyhdr}
\usepackage{titlesec}
\usepackage{tcolorbox}
% `enhanced` y `borderline west` viven en skins; `breakable` en su propia
% libreria. Sin declararlas, tcolorbox ignora esas claves en silencio y los
% recuadros pierden la barra lateral que los distingue.
\tcbuselibrary{skins,breakable}
\usepackage{booktabs}
\usepackage[hidelinks]{hyperref}

% ── Paleta ──────────────────────────────────────────────────────────────────
% Tomada de las variables de src/style.css. La interfaz es enteramente
% acromatica —todos sus colores son oklch(L 0 0)— asi que el manual tampoco
% introduce color: solo grises, con el texto como unico elemento de contraste.
\definecolor{fg}{HTML}{0A0A0A}        % --foreground
\definecolor{fgsoft}{HTML}{171717}    % --primary
\definecolor{muted}{HTML}{737373}     % --muted-foreground
\definecolor{surface}{HTML}{F5F5F5}   % --muted / --secondary
\definecolor{surfacealt}{HTML}{FAFAFA}% --sidebar
\definecolor{line}{HTML}{E5E5E5}      % --border

\titleformat{\section}{\LARGE\bfseries\color{fg}}{\thesection}{0.7em}{}
\titlespacing*{\section}{0pt}{1.6em}{0.7em}
\titleformat{\subsection}{\large\bfseries\color{fgsoft}}{\thesubsection}{0.6em}{}
\titlespacing*{\subsection}{0pt}{1.3em}{0.4em}
\titleformat{\subsubsection}[runin]{\bfseries\color{fgsoft}}{}{0pt}{}[.]

\pagestyle{fancy}
\fancyhf{}
\lhead{\footnotesize\color{muted}Manual de usuario}
\rhead{\footnotesize\color{muted}\thepage}
\renewcommand{\headrulewidth}{0.4pt}
\renewcommand{\headrule}{\color{line}\hrule height 0.4pt}

\setlength{\parindent}{0pt}
\setlength{\parskip}{0.62em}

% Enlaces y referencias en el mismo gris del cuerpo: sin subrayados ni azules.
\hypersetup{colorlinks=false}

% ── Recuadros ───────────────────────────────────────────────────────────────
% Se distinguen por una barra lateral y el peso del titulo, no por color.
\newtcolorbox{nota}[1][Nota]{
  enhanced, breakable, colback=surfacealt, colframe=surfacealt,
  borderline west={2pt}{0pt}{line},
  coltitle=fg, fonttitle=\bfseries\footnotesize\sffamily,
  title=#1, boxrule=0pt, arc=0pt, outer arc=0pt,
  left=10pt, right=10pt, top=7pt, bottom=7pt, toptitle=1pt, bottomtitle=3pt}

\newtcolorbox{aviso}[1][Importante]{
  enhanced, breakable, colback=surface, colframe=surface,
  borderline west={2pt}{0pt}{fg},
  coltitle=fg, fonttitle=\bfseries\footnotesize\sffamily,
  title=#1, boxrule=0pt, arc=0pt, outer arc=0pt,
  left=10pt, right=10pt, top=7pt, bottom=7pt, toptitle=1pt, bottomtitle=3pt}

% ── Captura de pantalla con marcador de posicion ────────────────────────────
% Uso: \captura{nombre-archivo}{Pie de la figura}
% Si manual/capturas/<nombre>.png existe, la inserta. Si no, dibuja el hueco.
\newcommand{\captura}[2]{%
  \begin{figure}[htbp]
  \centering
  \IfFileExists{capturas/#1.png}{%
    % Ancho completo, pero acotado en alto: una captura de ventana completa
    % puede ser muy alta y desbordaria la pagina. keepaspectratio evita que se
    % deforme cuando manda la restriccion de altura.
    {\color{line}\fboxrule=0.6pt\fboxsep=0pt\fbox{%
      \includegraphics[width=0.94\textwidth, height=0.62\textheight, keepaspectratio]{capturas/#1.png}}}%
  }{%
    \fcolorbox{line}{surfacealt}{%
      \begin{minipage}[c][3.1cm][c]{0.93\textwidth}
        \centering
        {\footnotesize\bfseries\color{muted}CAPTURA PENDIENTE}\\[5pt]
        {\ttfamily\footnotesize\color{fgsoft} capturas/#1.png}\\[4pt]
        {\footnotesize\color{muted} #2}
      \end{minipage}}%
  }
  \caption{#2}
  \label{fig:#1}
  \end{figure}%
}

% Pies de figura discretos
\usepackage{caption}
\captionsetup{font={small,color=muted}, labelfont={bf,color=muted}, skip=6pt}

 desde la
% carpeta de cada manual.
%
% La paleta sale de las variables de frontend/src/style.css: la interfaz es
% enteramente acromatica, y los manuales tampoco introducen color.
% ================================================================================
\documentclass[11pt, letterpaper]{article}

\usepackage[spanish,mexico]{babel}
\usepackage[utf8]{inputenc}
\usepackage[T1]{fontenc}

% Tipografia: Inter, la mas cercana en TeX Live a la Geist que usa la
% aplicacion. Sin serifas en todo el documento, igual que la interfaz.
\usepackage[sfdefault]{inter}
\usepackage[scaled=0.86]{beramono}

\usepackage[top=2.4cm, bottom=2.4cm, left=2.6cm, right=2.6cm]{geometry}
\usepackage{graphicx}
\usepackage{xcolor}
\usepackage{enumitem}
\usepackage{fancyhdr}
\usepackage{titlesec}
\usepackage{tcolorbox}
% `enhanced` y `borderline west` viven en skins; `breakable` en su propia
% libreria. Sin declararlas, tcolorbox ignora esas claves en silencio y los
% recuadros pierden la barra lateral que los distingue.
\tcbuselibrary{skins,breakable}
\usepackage{booktabs}
\usepackage[hidelinks]{hyperref}

% ── Paleta ──────────────────────────────────────────────────────────────────
% Tomada de las variables de src/style.css. La interfaz es enteramente
% acromatica —todos sus colores son oklch(L 0 0)— asi que el manual tampoco
% introduce color: solo grises, con el texto como unico elemento de contraste.
\definecolor{fg}{HTML}{0A0A0A}        % --foreground
\definecolor{fgsoft}{HTML}{171717}    % --primary
\definecolor{muted}{HTML}{737373}     % --muted-foreground
\definecolor{surface}{HTML}{F5F5F5}   % --muted / --secondary
\definecolor{surfacealt}{HTML}{FAFAFA}% --sidebar
\definecolor{line}{HTML}{E5E5E5}      % --border

\titleformat{\section}{\LARGE\bfseries\color{fg}}{\thesection}{0.7em}{}
\titlespacing*{\section}{0pt}{1.6em}{0.7em}
\titleformat{\subsection}{\large\bfseries\color{fgsoft}}{\thesubsection}{0.6em}{}
\titlespacing*{\subsection}{0pt}{1.3em}{0.4em}
\titleformat{\subsubsection}[runin]{\bfseries\color{fgsoft}}{}{0pt}{}[.]

\pagestyle{fancy}
\fancyhf{}
\lhead{\footnotesize\color{muted}Manual de usuario}
\rhead{\footnotesize\color{muted}\thepage}
\renewcommand{\headrulewidth}{0.4pt}
\renewcommand{\headrule}{\color{line}\hrule height 0.4pt}

\setlength{\parindent}{0pt}
\setlength{\parskip}{0.62em}

% Enlaces y referencias en el mismo gris del cuerpo: sin subrayados ni azules.
\hypersetup{colorlinks=false}

% ── Recuadros ───────────────────────────────────────────────────────────────
% Se distinguen por una barra lateral y el peso del titulo, no por color.
\newtcolorbox{nota}[1][Nota]{
  enhanced, breakable, colback=surfacealt, colframe=surfacealt,
  borderline west={2pt}{0pt}{line},
  coltitle=fg, fonttitle=\bfseries\footnotesize\sffamily,
  title=#1, boxrule=0pt, arc=0pt, outer arc=0pt,
  left=10pt, right=10pt, top=7pt, bottom=7pt, toptitle=1pt, bottomtitle=3pt}

\newtcolorbox{aviso}[1][Importante]{
  enhanced, breakable, colback=surface, colframe=surface,
  borderline west={2pt}{0pt}{fg},
  coltitle=fg, fonttitle=\bfseries\footnotesize\sffamily,
  title=#1, boxrule=0pt, arc=0pt, outer arc=0pt,
  left=10pt, right=10pt, top=7pt, bottom=7pt, toptitle=1pt, bottomtitle=3pt}

% ── Captura de pantalla con marcador de posicion ────────────────────────────
% Uso: \captura{nombre-archivo}{Pie de la figura}
% Si manual/capturas/<nombre>.png existe, la inserta. Si no, dibuja el hueco.
\newcommand{\captura}[2]{%
  \begin{figure}[htbp]
  \centering
  \IfFileExists{capturas/#1.png}{%
    % Ancho completo, pero acotado en alto: una captura de ventana completa
    % puede ser muy alta y desbordaria la pagina. keepaspectratio evita que se
    % deforme cuando manda la restriccion de altura.
    {\color{line}\fboxrule=0.6pt\fboxsep=0pt\fbox{%
      \includegraphics[width=0.94\textwidth, height=0.62\textheight, keepaspectratio]{capturas/#1.png}}}%
  }{%
    \fcolorbox{line}{surfacealt}{%
      \begin{minipage}[c][3.1cm][c]{0.93\textwidth}
        \centering
        {\footnotesize\bfseries\color{muted}CAPTURA PENDIENTE}\\[5pt]
        {\ttfamily\footnotesize\color{fgsoft} capturas/#1.png}\\[4pt]
        {\footnotesize\color{muted} #2}
      \end{minipage}}%
  }
  \caption{#2}
  \label{fig:#1}
  \end{figure}%
}

% Pies de figura discretos
\usepackage{caption}
\captionsetup{font={small,color=muted}, labelfont={bf,color=muted}, skip=6pt}

